\section{Lecture 1}
\subsection{Double Integrals Intro}

\begin{definition}{Double Integrals}{def311}
    The double summation of multivariable integrals is defined as:
    \begin{gather*}
        \sum_{i=1}^{\infty}\sum_{i=1}^{\infty}\Delta x\Delta y f(x_i,y_i) 
        =
        \iint_R f(x,y) DA
    \end{gather*}
    Where R is the independent x and y ranges or $\Delta x$ and $\Delta y$ and A is the area defined by those ranges.
    R does not have to be rectangular and can be defined by the area of functions.
\end{definition}

Double Integrals can be easily defined as twice iterated integrals.

\begin{lenma}{Simple dA definition under the x and y coordinate system}{lem311}
    \begin{gather*}
        dA = dxdy \quad \text{OR} \quad dA = dydx
    \end{gather*}
\end{lenma}

\begin{definition}{Volume of a box with height c}{def312}
    \begin{gather*}
        \iint_R c DA = c\iint_r DA = cA(R)
    \end{gather*}
\end{definition}

\begin{theorem}{Fubini's Theorem}{the311}
    Let $a \le x \le b$ and $g_1(x) \le y \le g_2(x)$ \\
    Let $c \le x \le d$ and $h_1(y) \le x \le h_2(y)$ \\
    Then for a continuous function $f(x,y)$ or R
    \begin{align*}
        \iint_R f(x,y) dA \qpand{=} \int_a^b\left[\int_{g_1(x)}^{g_2(x)} f(x,y)dy\right]dx \\
        \qpand{=} \int_c^d\left[\int_{h_1(y)}^{h_2(y)} f(x,y)dy\right]dx
    \end{align*}
    Alike partial derivatives, we can compare this to a partial integral.
\end{theorem}

\begin{example}{Example with figure given}{exam311}
    \begin{figure}[H]
        \def\svgwidth{\columnwidth}
        \import{Unit3/Lecture1/}{drawing1.pdf_tex}
        % \incfig{Unit3/Lecture1/example.pdf}
    \end{figure}
    \begin{align}
        \iint_R f(x,y) dA \qpand{=} \int_0^1\left[\int_{x^2}^{x} f(x,y)dy\right]dx \\
        \qpand{=} \int_0^1\left[\int_{y}^{\sqrt{y}} f(x,y)dy\right]dx
    \end{align}
\end{example}

\begin{definition}{Average Area}{def313}
    \begin{align*}
        \text{Average Area } = \frac{1}{A(R)}\iint_{R} f(x,y)\,dA
    \end{align*}
\end{definition}

\subsection{Polar Coordinates}
\begin{definitionside}{Polar Coordinates}{def314}
    \begin{figure}[H]
        \def\svgwidth{\columnwidth}
        \import{Unit3/Lecture1/}{drawing2test.pdf_tex}
    \end{figure}
    \tcblower
    \begin{align*}
        x = r\cos\theta \quad & \quad x^2+y^2=r^2 \\
        y = r\sin\theta \quad & \quad \frac{y}x = \tan\theta
    \end{align*}
\end{definitionside}



% \includepdf[width=100pt]{Unit3/Lecture1/example.pdf}Unit3/Lecture1/example